% Original source preamble and source-level macro declarations.
% This archival file is not the blueprint driver preamble.

\documentclass[a4paper,12pt]{article} %{article}
\title{Pixton's formula and Abel-Jacobi theory \\on the Picard stack}
%test_edit
%PACKAGES
%%%%%%%%%%%%%%%%%%%%%%%%%%%%%%%%%%%%%%%%%%%%%%%%
\usepackage{amsmath, amssymb, mathrsfs, amsthm, shorttoc, geometry, stmaryrd}% amsfonts, url,, tikz, wrapfig, newclude, }
\usepackage{mathtools} % for \coloneqq (and presumably other stuff!)
\usepackage{hyperref}
\usepackage{xcolor}
\hypersetup{
colorlinks,
linkcolor={black},
citecolor={blue!50!black},
urlcolor={blue!80!black}
}

\newcommand{\MWvone}[1]{}%{#1}
\newcommand{\MWref}[1]{#1}

\newcommand{\mat}[1]{\begin{bmatrix}#1\end{bmatrix}}

\usepackage[all]{xy}
%\usepackage[top=1.5cm, bottom=1.5cm]{geometry}
%\usepackage[left=2cm,right=2cm,top=3cm,bottom=3cm]{geometry}

\usepackage{tabularx}
\usepackage{longtable}
%\numberwithin{equation}{subsection}

%Enumitem
\usepackage{enumitem}
% eg \begin{enumerate}[label={\ref{definition:quasisplit}.\arabic*}]
% eg. \begin{enumerate}[label={1.\arabic*}]

%CLEVEREF STUFF	
\usepackage[capitalise]{cleveref}
%\usepackage{nameref}
\crefformat{equation}{(#2#1#3)}
\let\oref\ref
%\AtBeginDocument{\renewcommand{\ref}[1]{\cref{#1}}}

\usepackage{pgf,tikz}
\usetikzlibrary{matrix, calc, arrows}
%\usetikzlibrary{cd}
\DeclareMathOperator{\coker}{coker}

\usepackage{tikz-cd} %this is the old version, but it also the one the arXiv supports at the moment (2016)

%TODONOTES
%%%%%%%%
%\usepackage{todonotes}
%\setlength{\marginparwidth}{3cm}
%\reversemarginpar

% Double and triple arrows
\makeatletter
\newcommand*{\doublerightarrow}[2]{\mathrel{
  \settowidth{\@tempdima}{$\scriptstyle#1$}
  \settowidth{\@tempdimb}{$\scriptstyle#2$}
  \ifdim\@tempdimb>\@tempdima \@tempdima=\@tempdimb\fi
  \mathop{\vcenter{
    \offinterlineskip\ialign{\hbox to\dimexpr\@tempdima+1em{##}\cr
    \rightarrowfill\cr\noalign{\kern.5ex}
    \rightarrowfill\cr}}}\limits^{\!#1}_{\!#2}}}
\newcommand*{\triplerightarrow}[1]{\mathrel{
  \settowidth{\@tempdima}{$\scriptstyle#1$}
  \mathop{\vcenter{
    \offinterlineskip\ialign{\hbox to\dimexpr\@tempdima+1em{##}\cr
    \rightarrowfill\cr\noalign{\kern.5ex}
    \rightarrowfill\cr\noalign{\kern.5ex}
    \rightarrowfill\cr}}}\limits^{\!#1}}}
\makeatother
%eg
%$A\doublerightarrow{a}{bcdefgh}B$
%$A\triplerightarrow{d_0,d_1,d_2}B$

%MATH OPERATORS
\newcommand{\on}[1]{\operatorname{#1}}
\newcommand{\bb}[1]{{\mathbb{#1}}}
\newcommand{\cl}[1]{{\mathscr{#1}}}
\newcommand{\ca}[1]{{\mathcal{#1}}}
\newcommand{\bd}[1]{{\mathbf{#1}}}

\newcommand{\ul}[1]{{\underline{#1}}}


%Rahul abbreviations:
%%%%%%%%%%%%%%%%%%%%%%%%%%%%%%%%%%%%%%%%%%%%%%%%

\def\D{\mathrm{D}}
\def\d{\mathrm{d}}
\def\tw{\mathrm{tw}}
\def\red{\mathrm{red}}
\def\val{\mathrm{val}}
\def\log{\mathrm{log}}
\def\ev{\mathrm {ev}}
\def\vir{\mathrm {vir}}
\def\ct{\mathrm{ct}}
\def\virt{\mathrm{vir}}
\def\P{\mathsf{P}}
\def\PP{\mathbb{P}}
\def\DR{\mathsf{DR}}
\def\mathsff{{\mathcal{L}}}
\def\CP{{{\mathbb {CP}}}}
\def\oC{\overline{\mathcal{C}}}
\def\ooC{{\mathcal{C}}}
\def\cL{{\mathcal{L}}}
\def\cC{{\mathcal{C}}}
\def\cO{\mathcal{O}}
\def\oM{\overline{\mathcal{M}}}
\def\cM{{\mathcal{M}}}
\def\C{{\mathcal{C}}}
\def\eps{\varepsilon}
\def\Z{\mathbb{Z}}
\def\N{\mathbb{N}}
\def\ZZ{\mathcal{Z}}
\def\ZZZ{\mathsf{Z}}
\def\C{\mathbb{C}}
\def\Q{\mathbb{Q}}
%\def\d{\partial}
\def\qed{{\hfill $\Diamond$}}
\def\cS{{\mathcal S}}
\def\cg{{\mathcal G}}
\def\cR{{\mathcal R}}
\def\Aut{{\rm Aut}}
\def\F{{\mathsf{F}}}
\def\E{\mathrm{E}}
\def\n{\mathrm{n}}
\def\L{\mathrm{L}}
\def\V{\mathrm{V}}
\def\H{\mathrm{H}}
\def\HH{\mathcal{H}}
\def\g{\mathrm{g}}
\def\G{\mathsf{G}}
\def\rarr{\rightarrow}
\def\D{\mathsf{D}}
\def\tP{\widetilde{\mathsf{P}}}
\def\Klog{\omega_{\rm log}}
\def\rmW{\mathsf{W}}
\def\ff{\mathsf{Q}}
\def\fff{\widehat{\mathsf{Q}}}
\def\nice{\displaystyle}
\def\tC{\widetilde{\mathsf{C}}}
\def\ccoarseC{{\mathsf{C}}}
\def\coarseC{\overline{\mathsf{C}}}
\def\coarseL{\mathsf{L}}
\def\tL{\widetilde{{\mathcal L}}}
\def\tS{\widetilde{S}}
\def\tZ{\widetilde{Z}}
\def\oMM{\overline{\mathfrak{M}}}
\def\ooMM{{\mathfrak{M}}}
\def\oCC{\overline{\mathfrak{C}}}
\def\ooCC{{\mathfrak{C}}}
\def\oPP{\overline{\mathfrak{P}}}
\def\Pic{\mathcal{P}ic \,}
\def\ch{{\rm ch}}
\def\td{{\rm td}} 

\def\v{\mathsf{v}}
\def\h{\mathsf{h}}
\def\hH{\mathsf{h}}

%\def\bb{\mathsf{m}}
\def\cc{\ell}

\DeclareMathAlphabet{\mymathbb}{U}{BOONDOX-ds}{m}{n}
\def\zero{\mymathbb{0}}
\def\one{\mymathbb{1}}
\def\A{\mathcal{A}}

%SHEAF VERSIONS OF COMMON OPERATORS
%%%%%%%%%%%%%%%%%%%%%%%%%%%%%%%%%%%%%%%%%%%%%%%%
\def\sheaflie{\mathcal{L}ie}
\def\sheafhom{\mathcal{H}om}

%BRACKETS AND PAIRINGS
%%%%%%%%%%%%%%%%%%%%%%%%%%%%%%%%%%%%%%%%%%%%%%%%
\def\lleft{\left<\!\left<}%{\langle\langle}
\def\rright{\right>\!\right>}%{\rangle\rangle}
\newcommand{\Span}[1]{\left<#1\right>}
\newcommand{\SSpan}[1]{\lleft#1\rright}
\newcommand{\abs}[1]{\lvert#1\rvert}
\newcommand{\aabs}[1]{\lvert\lvert#1\rvert\rvert}

%SUB- AND SUPERSCRIPTS
%%%%%%%%%%%%%%%%%%%%%%%%%%%%%%%%%%%%%%%%%%%%%%%%
\newcommand{\algcl}{^{\textrm{alg}}}
\newcommand{\sepcl}{{^{\textrm{sep}}}}
\newcommand{\urcl}{{^{\textrm{ur}}}}
\newcommand{\ssm}{_{\smooth}}

%FONT ENCODING AND EXOTIC LETTERS
%%%%%%%%%%%%%%%%%%%%%%%%%%%%%%%%%%%%%%%%%%%%%%%%
%\usepackage[OT2,OT1]{fontenc}
%\DeclareSymbolFont{cyrletters}{OT2}{wncyr}{m}{n}
%\DeclareMathSymbol{\Sha}{\mathalpha}{cyrletters}{"58}

\newcommand{\ra}{\rightarrow} %depreciated; use \to
\newcommand{\lra}{\longrightarrow}
\newcommand{\hra}{\hookrightarrow}
\newcommand{\sub}{\subseteq}
\newcommand{\tra}{\rightarrowtail}
\newcommand{\iso}{\stackrel{\sim}{\lra}}


%THEOREMS
%%%%%%%%%%%%%%%%%%%%%%%%%%%%%%%%%%%%%%%%%%%%%%%%
\theoremstyle{definition}
\newtheorem{definition}{Definition} %[section]
\newtheorem{condition}[definition]{Condition}
\newtheorem{problem}[definition]{Problem}
\newtheorem{situation}[definition]{Situation}
\newtheorem{construction}[definition]{Construction}


\newtheorem{remark}[definition]{Remark}
\newtheorem{aside}[definition]{Aside}
\newtheorem{example}[definition]{Example}

\theoremstyle{plain}% default
\newtheorem{conjecture}[definition]{Conjecture}
\newtheorem*{conjectureA}{Conjecture A}
\newtheorem{proposition}[definition]{Proposition}
\newtheorem{lemma}[definition]{Lemma}
\newtheorem{theorem}[definition]{Theorem}
\newtheorem*{theorem*}{Theorem}
\newtheorem{corollary}[definition]{Corollary}
\newtheorem{claim}[definition]{Claim}
\newtheorem{inclaim}{Claim}[definition]

%\theoremstyle{remark}


%MISC
%%%%%%%%%%%%%%%%%%%%%%%%%%%%%%%%%%%%%%%%%%%%%%%%
\def\tam{c_{\tau}} %Tamagawa number
\def\defeq{\stackrel{\text{\tiny def}}{=}}%depreciated: use coloneqq
\newcommand{\linequiv}{\sim_{\text{\tiny lin}}}
\newcommand{\MbarX}{\Mbar_{g,n}(X, \beta)}

%% Some convolutions to redefine \phi as \varphi while storing the original \phi in \phiorig
\usepackage{letltxmacro}
\LetLtxMacro{\phiorig}{\phi}
% \newcommand{\phiorig}{\phi}
\renewcommand{\phi}{\varphi}
\newcommand{\aj}{{\scriptstyle \mathsf{AJ}}}
\newcommand{\vdim}{\mathrm{vdim}}


\newcommand{\cecH}{\check{\on H}}

%\newcounter{temp} 
%\newcommand\interitem[1]{\setcounter{temp}{\value{enumi}}\end{enumerate}{#1}\begin{enumerate}\setcounter{enumi}{\value{temp}}}
%Then use \interitem{text goes here} between two items in a top-level enumerate environment.  It doesn't work inside nested enumerate environments or if you forget to have another item after your interitem text.

\newcommand{\loz}{\lozenge}
\newcommand{\bloz}{\blacklozenge}
\newcommand{\barsigma}{\bar{\sigma}}

\newcommand{\Picabs}{\mathfrak{Pic}}
\newcommand{\Picrel}{\mathfrak{Pic}^{\mathrm{rel}}}
\newcommand{\Chow}{\mathsf{CH}}
\newcommand{\CHop}{\Chow_{\mathsf{op}}}

\newcommand{\tanmap}{\psi}

\newcommand{\invnrI}{I}
\newcommand{\invnrn}{II}
\newcommand{\invnrA}{III}
\newcommand{\invnrII}{IV}
\newcommand{\invnrIII}{V}
\newcommand{\invnrIV}{VI}

% \newcommand{\jocomment}[1]{\textcolor{red}{Jo: #1}}
%\newcommand{\jocomment}[1]{{\color{red}Jo: #1}}
 \newcommand{\jocomment}[1]{}
%\newcommand{\Dcomment}[1]{{\color{blue}D: #1}}
 \newcommand{\Dcomment}[1]{}
%\newcommand{\Rcomment}[1]{{\textbf{\color{red}RP: #1}}}
%\newcommand{\Rosacomment}[1]{{\color{violet}Rosa: #1}}
%\newcommand{\Ycomment}[1]{{\color{brown}Y: #1}}
%AUTHOR DETAILS
%%%%%%%%%%%%%%%%%%%%%%%%%%%%%%%%%%%%%%%%%%%%%%%%
\author{Y. Bae, D. Holmes, R. Pandharipande, J. Schmitt, R. Schwarz}
\date{May 2021}



%COMMENTS
%%%%%%%%%%%%%%%%%%%%%%%%%%%%%%%%%%%%%%%%%%%%%%%%
\newcounter{nootje}
\setcounter{nootje}{1}
\newcommand\todo[1]{[*\thenootje]\marginpar{\tiny\begin{minipage}
{20mm}\begin{flushleft}\thenootje : 
#1\end{flushleft}\end{minipage}}\addtocounter{nootje}{1}}
%\renewcommand{\todo}[1]{}


%Longer versions of proofs
%%%%%%%%%%%%%%%%%%%%%%%%%%%%%%%%%%%%%%%%%%%%%%%%
% \newcommand{\expanded}[1]{#1}
% \renewcommand{\expanded}[1]{}

%SHORTCUTS FOR COMMON COMMANDS
%%%%%%%%%%%%%%%%%%%%%%%%%%%%%%%%%%%%%%%%%%%%%%%%
\newcommand{\beq}{\begin{equation}}%depreciated
\newcommand{\eeq}{\end{equation}}%depreciated
\newcommand{\beqs}{\begin{equation*}}%depreciated
\newcommand{\eeqs}{\end{equation*}}%depreciated

\renewcommand{\k}{k}

\newcommand{\DRop}{\mathsf{DR}^{\mathsf{op}}}

% For using \subset in tikzcd
\tikzset{
  symbol/.style={
    draw=none,
    every to/.append style={
      edge node={node [sloped, allow upside down, auto=false]{$#1$}}}
  }
}
%For rotating labels in tikzcd - added by Rosa 01-04-2020
\tikzset{
    labl/.style={anchor=south, rotate=-90, inner sep=.5mm}
}


